\documentclass[11pt,a4paper]{article}

% Essential packages
\usepackage[utf8]{inputenc}
\usepackage[T1]{fontenc}
\usepackage{amsmath,amsthm,amssymb}
\usepackage{hyperref}
\usepackage{cleveref}

% Theorem environments
\newtheorem{theorem}{Theorem}[section]
\newtheorem{lemma}[theorem]{Lemma}
\newtheorem{proposition}[theorem]{Proposition}
\newtheorem{corollary}[theorem]{Corollary}
\theoremstyle{definition}
\newtheorem{definition}[theorem]{Definition}
\newtheorem{example}[theorem]{Example}
\theoremstyle{remark}
\newtheorem{remark}[theorem]{Remark}

% Open question environment
\newenvironment{openquestion}{\par\textbf{Open Question.}\itshape}{\par}

\title{Fiber Ambiguity as a Hiding Heuristic:\\
Linear Maps vs.\ a 2-to-1 Covering \(T^2\to K\)}
\author{AI Notes}
\date{2026-01-15}

\begin{document}
\maketitle

\begin{abstract}
This note records a ``fiber ambiguity'' heuristic: when an observer learns only an output value of a map \(f\), the set of consistent inputs \(f^{-1}(s)\) (the fiber) may contain multiple candidates, potentially enabling information hiding. A linear reference example \(f:\mathbb{R}^2\to\mathbb{R}\) is contrasted with a dimension-preserving but many-to-one map arising from topology: a 2-to-1 covering map from the torus \(T^2\) onto the Klein bottle \(K\). The purpose is not to claim cryptographic security from topology alone, but to isolate the mechanism ``finite fibers without dimension drop'' and to provide an explicit, programmatic description of the covering map.
\end{abstract}

\section{Motivation and scope}
\label{sec:motivation}

Suppose two parties (``Bob'' and ``Charlie'') hold private values \(x\) and \(y\), and an observer (``Alice'') learns only an output
\[
s = f(x,y)
\]
for some publicly known function \(f\). A first heuristic for ``hiding'' is that if \(f\) is not injective, then for a given \(s\) there may be multiple \((x,y)\) compatible with the observation; in geometric language, the fiber \(f^{-1}(s)\) may contain more than one point.

This note makes two points.
\begin{itemize}
  \item The size/dimension of fibers is a \emph{geometric} way to see non-injectivity.
  \item ``Hiding'' is ultimately \emph{distributional}: one must analyze what remains uncertain about \((x,y)\) after observing \(s\), e.g. via conditional entropy. Fiber size alone is not a complete security argument.
\end{itemize}

\section{A reference linear example}
\label{sec:linear}

\subsection{Fibers of a nonzero linear functional}
\label{sec:linear-fibers}

Fix \((a,c)\in \mathbb{R}^2\setminus\{(0,0)\}\) and define
\[
f:\mathbb{R}^2\to\mathbb{R},\qquad f(x,y)=ax+cy.
\]

\begin{proposition}[Affine-line fibers]
\label{prop:affine-line-fibers}
For each \(s\in \mathbb{R}\), the fiber \(f^{-1}(s)=\{(x,y)\in\mathbb{R}^2: ax+cy=s\}\) is an affine line in \(\mathbb{R}^2\). In particular, it is infinite and 1-dimensional.
\end{proposition}

\begin{proof}
Since \((a,c)\neq(0,0)\), the equation \(ax+cy=s\) is a single nontrivial linear constraint in \((x,y)\). Its solution set is a translate of the kernel \(\ker(f)=\{(x,y): ax+cy=0\}\), which is a 1-dimensional linear subspace of \(\mathbb{R}^2\). Thus \(f^{-1}(s)\) is an affine line.
\end{proof}

\subsection{A finite-domain ``masking'' variant (distributional)}
\label{sec:linear-masking}

The previous subsection is purely geometric. To connect to ``hiding'', it is useful to state one clean finite-domain fact.

\begin{proposition}[Perfect hiding by an additive mask]
\label{prop:perfect-hiding}
Let \(N\ge 2\). Let \(X,Y\) be random variables on \(\mathbb{Z}_N\) such that \(Y\) is uniform on \(\mathbb{Z}_N\) and independent of \(X\). Fix \(a,c\in\mathbb{Z}_N\) with \(c\) invertible in \(\mathbb{Z}_N\) (equivalently \(\gcd(c,N)=1\)). Define
\[
S \equiv aX + cY \pmod N.
\]
Then \(S\) is uniform on \(\mathbb{Z}_N\) and independent of \(X\). Equivalently, \(I(X;S)=0\).
\end{proposition}

\begin{proof}
Fix \(x\in\mathbb{Z}_N\). Conditional on \(X=x\), one has \(S\equiv ax + cY\ (\mathrm{mod}\ N)\). Since \(c\) is invertible, the map \(y\mapsto ax+cy\) is a bijection of \(\mathbb{Z}_N\), hence it pushes forward the uniform distribution on \(Y\) to the uniform distribution on \(S\). Therefore \(\mathbb{P}(S=s\mid X=x)=1/N\) for all \(s\), which implies that \(S\) is uniform and independent of \(X\).
\end{proof}

\begin{remark}
\label{rem:fiber-vs-entropy}
In \cref{prop:perfect-hiding}, large fibers are not the primary driver: the key is the presence of an independent uniform ``mask'' \(cY\). Conversely, if \(Y\) is not close to uniform or is correlated with \(X\), then observing \(S\) can leak substantial information even if each fiber contains many points.
\end{remark}

\section{A 2-to-1 covering \(T^2\to K\): finite fibers without dimension drop}
\label{sec:cover}

\subsection{Modeling \(T^2\) and \(K\) by fundamental domains}
\label{sec:domains}

Write the torus as a quotient
\[
T^2 := \bigl(\mathbb{R}/\mathbb{Z}\bigr)\times\bigl(\mathbb{R}/\mathbb{Z}\bigr),
\]
so a point of \(T^2\) is represented by \((u,v)\in [0,1)\times[0,1)\) with arithmetic performed modulo \(1\) in each coordinate.

Write the Klein bottle as the quotient of \([0,1)\times[0,1)\) under the edge identifications
\[
(x,0)\sim(x,1),\qquad (0,y)\sim(1,1-y).
\]
This is a standard fundamental polygon description of the Klein bottle; for the present note, it can be treated as a definition of a space \(K\).

\subsection{An explicit 2-to-1 map}
\label{sec:explicit-map}

Define an involution \(\tau:T^2\to T^2\) by
\[
\tau(u,v) := \bigl(u+\tfrac12,\, 1-v\bigr) \quad (\mathrm{mod}\ 1).
\]
This map is fixed-point free: if \(\tau(u,v)=(u,v)\), then \(u+\tfrac12\equiv u\pmod 1\) (impossible) and \(1-v\equiv v\pmod 1\) (i.e.\ \(v\in\{0,\tfrac12\}\)), but the first condition already fails. Hence \(\tau\) is a free \(\mathbb{Z}_2\)-action on \(T^2\), and the quotient \(T^2/\langle\tau\rangle\) is a \(2\)-sheeted covering space.

Concretely, define a map \(p:T^2\to K\) as follows. Given \((u,v)\in[0,1)\times[0,1)\) representing a point of \(T^2\), first \emph{normalize} it by folding along \(u=\tfrac12\):
\[
(u',v') :=
\begin{cases}
(u,v) & \text{if } u\in[0,\tfrac12),\\
(u-\tfrac12,\, (1-v)\bmod 1) & \text{if } u\in[\tfrac12,1).
\end{cases}
\]
Then set
\[
p(u,v) := [\, (x,y) \,] \in K,\qquad (x,y):=(2u',\, v').
\]
Here \([ (x,y) ]\) denotes the equivalence class in \(K\) under the edge identifications in \cref{sec:domains}. The map \(p\) descends from the quotient by \(\tau\), and since \(\tau\) is a free involution, \(p:T^2\to K\) is a \(2\)-sheeted covering map.

\begin{proposition}[Two preimages per point (heuristic fiber statement)]
\label{prop:two-preimages}
For each \([ (x,y) ]\in K\) represented by \((x,y)\in[0,1)\times[0,1)\), the two points
\[
(u_1,v_1):=\Bigl(\tfrac{x}{2},\, y\Bigr),\qquad
(u_2,v_2):=\Bigl(\tfrac{x}{2}+\tfrac12,\, 1-y\Bigr)
\]
in \([0,1)\times[0,1)\) represent points of \(T^2\) satisfying \(p(u_1,v_1)=p(u_2,v_2)=[(x,y)]\). In particular, \(p\) is not injective.
\end{proposition}

\begin{proof}
By construction, \((u_1,v_1)\) lies in the first half \(u\in[0,\tfrac12)\), so normalization does not change it; hence \(p(u_1,v_1)=[(2u_1,v_1)]=[(x,y)]\).

For \((u_2,v_2)\), one has \(u_2\in[\tfrac12,1)\), so normalization sends \((u_2,v_2)\mapsto (u_2-\tfrac12,1-v_2)=(\tfrac{x}{2},y)\). Thus \(p(u_2,v_2)=[(x,y)]\) as well.
\end{proof}

\begin{remark}
\label{rem:covering-vs-dimension-drop}
The map \(p:T^2\to K\) has \(\dim(T^2)=\dim(K)=2\) but still has nontrivial fibers (here, generically two points). This contrasts with the linear map in \cref{sec:linear}, where \(\mathbb{R}^2\to\mathbb{R}\) forces positive-dimensional fibers for nondegenerate linear constraints.
\end{remark}

\subsection{Covering-map viewpoint: fiber ambiguity and unique lifting}
\label{sec:hiding-heuristic}

In a hypothetical protocol inspired by \cref{prop:two-preimages}, Bob and Charlie might hold \((u,v)\in T^2\) while Alice learns only \(p(u,v)\in K\). The fiber \(p^{-1}(p(u,v))\) contains (at least) two candidates related by the involution
\[
(u,v)\longmapsto \Bigl(u+\tfrac12,\, 1-v\Bigr)\quad (\mathrm{mod}\ 1).
\]

A complementary (and more ``dynamic'') viewpoint is available when \(p:E\to B\) is a covering map: paths and homotopies in \(B\) can be \emph{lifted uniquely once an initial lift is fixed}. This is the formalization of ``unrolling''.

\begin{definition}[Unique path lifting property]
\label{def:uplp}
A map \(p:E\to B\) has the \emph{unique path lifting property} if for every path \(\gamma:I\to B\) and every choice of \(\tilde e_0\in E\) with \(p(\tilde e_0)=\gamma(0)\), there exists a unique path \(\tilde\gamma:I\to E\) such that
\[
p\circ \tilde\gamma = \gamma,
\qquad
\tilde\gamma(0)=\tilde e_0.
\]
\end{definition}

\begin{remark}[Diagrammatic form]
\label{rem:uplp-diagram}
\Cref{def:uplp} is equivalently the statement that for \(i_0:\{0\}\hookrightarrow I\), every solid diagram admits a \emph{unique} dotted lift making it commute:
\[
\begin{array}{ccc}
\{0\} & \xrightarrow{\ \tilde f\ } & E \\
\Big\downarrow {i_0} & \nearrow_{\ \exists!\ \tilde \gamma\ } & \Big\downarrow {p} \\
I & \xrightarrow{\ \gamma\ } & B.
\end{array}
\]
Here \(\tilde f(*)=\tilde e_0\) and commutativity encodes \(p\circ \tilde\gamma=\gamma\) together with the initial condition \(\tilde\gamma(0)=\tilde e_0\).
\end{remark}

\begin{proposition}[Covering maps lift paths uniquely]
\label{prop:covering-uplp}
If \(p:E\to B\) is a covering map, then \(p\) has the unique path lifting property.
\end{proposition}

\begin{proof}[Proof sketch]
Cover \(\gamma(I)\subset B\) by evenly covered neighborhoods. By compactness, finitely many suffice; subdivide \(I\) so that each subinterval maps into one such neighborhood. On each subinterval, the lift is forced by continuity and the requirement that it stay in the same sheet. Uniqueness follows because distinct lifts would have to separate at some point, contradicting local homeomorphism of sheets.
\end{proof}

A related uniqueness statement holds for lifting homotopies \(H:X\times I\to B\): informally, if one fixes a lift of \(H(\cdot,0)\), then the lift of the entire homotopy is forced (componentwise) by the covering-map structure.

\paragraph{Protocol interpretation (dynamic viewpoint).}
In a protocol-flavored reading, one may regard the parties' \emph{joint secret state} as a time-indexed point \(\tilde\gamma(t)=(u(t),v(t))\in T^2\), while Alice observes only its image \(\gamma(t)=p(\tilde\gamma(t))\in K\). Then \cref{prop:covering-uplp} says that, once the parties' initial secret \(\tilde\gamma(0)\in p^{-1}(\gamma(0))\) is fixed (equivalently, once the initial ``sheet'' is fixed), the entire lifted trajectory \(\tilde\gamma\) is determined uniquely by the observed base trajectory \(\gamma\). Without knowledge of the initial sheet, Alice can at best narrow down the evolution to the (finite) set of lifts consistent with \(\gamma\) and the unknown initial choice.

\begin{remark}[Two possible ``unrollings'' in a \(2\)-sheeted cover]
\label{rem:two-unrollings}
For a \(2\)-sheeted cover \(p:E\to B\), observing only a base path \(\gamma\) does \emph{not} determine a unique lift in \(E\): one must also choose the initial point \(\tilde e_0\in p^{-1}(\gamma(0))\). Once \(\tilde e_0\) is fixed, the lift is unique by \cref{def:uplp,prop:covering-uplp}. Thus, in the \(T^2\to K\) picture, the basic ambiguity can be viewed as the discrete choice of ``sheet'' at time \(0\); after that choice, the evolution is determined.
\end{remark}

\begin{remark}[What fiber ambiguity does and does not imply]
\label{rem:not-security}
The presence of a 2-element fiber suggests a one-bit ambiguity. However, turning this into a security statement requires a distributional assumption: one must specify a distribution on \((U,V)\) and analyze the posterior distribution given \(p(U,V)\). In particular, it is possible for \(p(U,V)\) to leak essentially all of \(U,V\) except a single bit, depending on what Alice can observe and how \((U,V)\) is generated. Section~\ref{sec:info-theory} makes this precise in a finite discretization.
\end{remark}

\section{Information-theoretic analysis (finite discretization)}
\label{sec:info-theory}

This section analyzes hiding in a finite-alphabet model, in the same spirit as \cref{prop:perfect-hiding}. The use of a finite model is deliberate: for deterministic maps between continuous random variables, differential-entropy expressions for mutual information require care (indeed, for \(K=p(U,V)\) deterministic, the conditional differential entropy \(h(K\mid U,V)\) is not a finite number), whereas for finite alphabets the Shannon mutual information is unambiguous.

\subsection{A discrete torus and a discrete Klein bottle}
\label{sec:discrete-model}

Fix an integer \(n\ge 2\). Let
\[
\mathbb{T}_n := \mathbb{Z}_{2n}\times \mathbb{Z}_{n}
\]
be a discrete model of a torus (the first factor is taken of size \(2n\) to accommodate a ``fold at \(1/2\)''). Let
\[
\mathbb{K}_n := \mathbb{Z}_{n}\times \mathbb{Z}_{n}
\]
be a discrete model of the Klein bottle fundamental square.

\begin{definition}[Discrete 2-to-1 map]
\label{def:discrete-cover}
Define \(p_n:\mathbb{T}_n\to \mathbb{K}_n\) by
\[
p_n(u,v) :=
\begin{cases}
(u,\ v) & \text{if } u\in\{0,1,\dots,n-1\},\\
(u-n,\ -v) & \text{if } u\in\{n,n+1,\dots,2n-1\},
\end{cases}
\]
where \(-v\) is computed in \(\mathbb{Z}_n\).
\end{definition}

\begin{remark}[Reproducibility]
\label{rem:reproducibility}
The identities proved in \cref{prop:discrete-two-fiber,prop:hide-bit,prop:leakage-u} are verified by exhaustive enumeration in the accompanying script \texttt{discrete\_cover\_info.py}. Running \texttt{python3 discrete\_cover\_info.py} in the document directory prints \(H(\cdot)\) and \(I(\cdot;\cdot)\) values for several \(n\) alongside the predicted formulas.
\end{remark}

\begin{proposition}[Two-point fibers]
\label{prop:discrete-two-fiber}
For each \((x,y)\in\mathbb{K}_n\), the fiber \(p_n^{-1}(x,y)\subset \mathbb{T}_n\) consists of exactly two points:
\[
p_n^{-1}(x,y)=\{(x,y),\ (x+n,-y)\}.
\]
\end{proposition}

\begin{proof}
By inspection, \(p_n(x,y)=(x,y)\) (first case) and \(p_n(x+n,-y)=(x,y)\) (second case). Conversely, if \(p_n(u,v)=(x,y)\), then either \(u\in\{0,\dots,n-1\}\) and \((u,v)=(x,y)\), or \(u\in\{n,\dots,2n-1\}\) and then \((u-n,-v)=(x,y)\), i.e.\ \((u,v)=(x+n,-y)\).
\end{proof}

\subsection{What is hidden depends on what is declared ``the secret''}
\label{sec:what-is-secret}

Let \((U,V)\) be random variables uniformly distributed on \(\mathbb{T}_n\), and let Alice observe
\[
K := p_n(U,V)\in \mathbb{K}_n.
\]
Write the first coordinate \(U\in \mathbb{Z}_{2n}\) uniquely as
\[
U = T + nB,
\qquad
T\in\mathbb{Z}_n,\quad B\in\{0,1\},
\]
so that \(T\) records the class \(U\bmod n\) and \(B\) records which ``sheet'' (first half vs second half) the input is on. Note that under the uniform distribution on \(\mathbb{T}_n\), the triple \((B,T,V)\) is independent with \(B\sim \mathrm{Bernoulli}(1/2)\), \(T\sim \mathrm{Unif}(\mathbb{Z}_n)\), and \(V\sim \mathrm{Unif}(\mathbb{Z}_n)\).

\begin{proposition}[Perfect hiding of the sheet bit]
\label{prop:hide-bit}
In the above model, the observed output \(K\) is independent of the sheet bit \(B\). Equivalently,
\[
I(B;K)=0.
\]
\end{proposition}

\begin{proof}
By \cref{def:discrete-cover}, one has \(K=(T, W)\) where \(W=V\) if \(B=0\) and \(W=-V\) if \(B=1\). Since \(V\) is uniform on \(\mathbb{Z}_n\), so is \(-V\), hence \(W\) is uniform on \(\mathbb{Z}_n\) and independent of \(B\). Also \(T\) is independent of \(B\). Therefore \((T,W)\) is independent of \(B\), i.e.\ \(K\) is independent of \(B\).
\end{proof}

\begin{proposition}[One-bit ambiguity about \((U,V)\), substantial leakage about \(U\)]
\label{prop:leakage-u}
In the same model,
\[
H(U,V\mid K)=1,
\qquad
H(U\mid K)=1,
\]
and therefore
\[
I\bigl((U,V);K\bigr)=2\log n,
\qquad
I(U;K)=\log n.
\]
\end{proposition}

\begin{proof}
By \cref{prop:discrete-two-fiber}, each value \(k\in\mathbb{K}_n\) has exactly two preimages in \(\mathbb{T}_n\). Under the uniform distribution on \(\mathbb{T}_n\), these two preimages are equiprobable conditional on \(K=k\), hence \(H(U,V\mid K)=1\).

The same fiber description shows that conditional on \(K\), the coordinate \(U\in\mathbb{Z}_{2n}\) has exactly two equiprobable values (differing by \(n\)), hence \(H(U\mid K)=1\).

Finally, \(H(U,V)=\log(2n^2)=1+2\log n\) and \(H(U)=\log(2n)=1+\log n\). Thus
\[
I((U,V);K)=H(U,V)-H(U,V\mid K)=(1+2\log n)-1=2\log n,
\]
and
\[
I(U;K)=H(U)-H(U\mid K)=(1+\log n)-1=\log n.
\]
\end{proof}

\begin{remark}
\label{rem:interpretation}
\Cref{prop:hide-bit,prop:leakage-u} formalize the ``fiber-ish'' intuition precisely in a finite model. The map \(p_n\) hides at most one bit: it perfectly hides the \(\mathbb{Z}_2\) ``sheet'' variable \(B\), but it reveals the residue \(T=U\bmod n\) completely. In particular, if Bob's secret is the full \(U\), then \(K\) leaks \(\log n\) bits about it. If Bob's secret is only the sheet bit \(B\), then \(K\) is perfectly hiding under the stated distributional assumptions.
\end{remark}

\begin{example}[Worked values for \(T,B,U,V\) and the output \(K=p_n(U,V)\)]
\label{ex:table-t-b-u-v}
For a concrete illustration, take \(n=5\). Then \(\mathbb{T}_5=\mathbb{Z}_{10}\times\mathbb{Z}_5\) and \(\mathbb{K}_5=\mathbb{Z}_5\times\mathbb{Z}_5\). Every \(U\in\mathbb{Z}_{10}\) can be written uniquely as \(U=T+5B\) with \(T\in\mathbb{Z}_5\) and \(B\in\{0,1\}\). Moreover, \(-V\) denotes the additive inverse of \(V\) in \(\mathbb{Z}_5\).

\medskip
\noindent\textbf{Legend of symbols:}
\begin{itemize}
  \item \(U\in\mathbb{Z}_{10}\): Bob's \emph{private input} (the first coordinate of a point on the discrete torus \(\mathbb{T}_5\)).
  \item \(V\in\mathbb{Z}_5\): Charlie's \emph{private input} (the second coordinate of a point on the discrete torus \(\mathbb{T}_5\)).
  \item \((U,V)\in\mathbb{T}_5\): the \emph{joint private input} (a point on the discrete torus).
  \item \(T=U\bmod 5\): the \emph{residue class}, which is \textbf{revealed} to Alice via the first coordinate of \(K\).
  \item \(B\in\{0,1\}\): the \emph{sheet bit}, indicating which half of the torus the input lies in; this is \textbf{hidden} from Alice (under uniform \(V\)).
  \item \(K=p_5(U,V)\in\mathbb{K}_5\): the \emph{observed output} (a point on the discrete Klein bottle), observed by Alice.
\end{itemize}

\begin{center}
\begin{tabular}{r r r r r l}
\hline
\(U\) & \(V\) & \(T=U\bmod 5\) & \(B\) & \(-V \bmod 5\) & \(K=p_5(U,V)\) \\
 \footnotesize (torus) & \footnotesize (torus) & \footnotesize (revealed) & \footnotesize (hidden) & & \footnotesize (Klein bottle) \\
\hline
0 & 0 & 0 & 0 & 0 & \((0,0)\) \\
0 & 1 & 0 & 0 & 4 & \((0,1)\) \\
2 & 4 & 2 & 0 & 1 & \((2,4)\) \\
4 & 3 & 4 & 0 & 2 & \((4,3)\) \\
5 & 0 & 0 & 1 & 0 & \((0,0)\) \\
5 & 1 & 0 & 1 & 4 & \((0,4)\) \\
7 & 2 & 2 & 1 & 3 & \((2,3)\) \\
9 & 4 & 4 & 1 & 1 & \((4,1)\) \\
\hline
\end{tabular}
\end{center}

The definition \cref{def:discrete-cover} implies that \(K=(T,V)\) when \(B=0\) and \(K=(T,-V)\) when \(B=1\). In particular, the first coordinate of \(K\) reveals \(T=U\bmod 5\) exactly, while the sheet bit \(B\) is not determined by \(K\): for instance, \(p_5(0,1)=(0,1)\) and also \(p_5(5,4)=(0,1)\), so observing \((0,1)\) does not distinguish between \(B=0\) and \(B=1\) in this case.
\end{example}

\subsection{Biased inputs: when the sheet bit remains hidden (and when it leaks)}
\label{sec:biased-u}

The results above assume that \((U,V)\) is uniform on \(\mathbb{T}_n\). In potential ``protocol'' interpretations, however, the distribution of \((U,V)\) may be biased. This subsection records the simplest information-theoretic conditions under which the sheet bit \(B\) remains hidden from the observation \(K=p_n(U,V)\).

\begin{proposition}[General distribution: what determines the posterior on the sheet bit]
\label{prop:posterior-b-general}
Let \(n\ge 2\). Let \((U,V)\) be \emph{any} random variables on \(\mathbb{Z}_{2n}\times \mathbb{Z}_n\). Write \(U=T+nB\) with \(T\in\mathbb{Z}_n\), \(B\in\{0,1\}\), and let \(K=p_n(U,V)=(T,W)\) where \(W=V\) if \(B=0\) and \(W=-V\) if \(B=1\). Then for any \((t,w)\in \mathbb{Z}_n\times\mathbb{Z}_n\) with \(\mathbb{P}(K=(t,w))>0\),
\[
\mathbb{P}(B=1\mid K=(t,w))
=
\frac{\mathbb{P}(B=1,\,T=t,\,V=-w)}{\mathbb{P}(B=0,\,T=t,\,V=w)+\mathbb{P}(B=1,\,T=t,\,V=-w)}.
\]
\end{proposition}

\begin{proof}
By definition, the event \(\{K=(t,w),\,B=0\}\) is the same as \(\{B=0,\,T=t,\,V=w\}\), while \(\{K=(t,w),\,B=1\}\) is the same as \(\{B=1,\,T=t,\,V=-w\}\). Bayes' rule then yields the displayed expression.
\end{proof}

\begin{remark}[Two distinct failure modes]
\label{rem:two-failure-modes}
\Cref{prop:posterior-b-general} isolates two independent mechanisms by which \(B\) can become learnable from \(K\).
\begin{itemize}
  \item If \(B\) is correlated with \(T\) (equivalently, the ``sheet'' is biased \emph{depending on} the residue class \(U\bmod n\)), then the \emph{first} coordinate of \(K\), which equals \(T\), can leak information about \(B\).
  \item If the conditional distribution of \(V\) given \((B,T)\) is not symmetric under negation (so that \(V\) and \(-V\) are distinguishable), then the \emph{second} coordinate of \(K\) can leak information about \(B\).
\end{itemize}
\end{remark}

\begin{proposition}[A simple sufficient condition for perfect hiding of \(B\)]
\label{prop:biased-u-sufficient}
Assume that \(V\) is independent of \((B,T)\) and that \(V\) is symmetric under negation, i.e.\ \(V\overset{d}{=}-V\) as \(\mathbb{Z}_n\)-valued random variables. Then \(W\) is independent of \(B\) (indeed \(W\overset{d}{=}V\)), and consequently
\[
I(B;K)=I(B;T).
\]
In particular, under these assumptions, the sheet bit \(B\) is perfectly hidden from \(K\) (i.e.\ \(I(B;K)=0\)) if and only if \(B\) is independent of \(T\).
\end{proposition}

\begin{proof}
Under the stated assumptions, the conditional law of \(W\) given \((B,T)\) does not depend on \(B\): if \(B=0\) then \(W=V\), and if \(B=1\) then \(W=-V\), but \(V\overset{d}{=}-V\). Hence \(I(B;W\mid T)=0\). Using the chain rule,
\[
I(B;K)=I(B;T,W)=I(B;T)+I(B;W\mid T)=I(B;T).
\]
The final claim follows immediately.
\end{proof}

\begin{remark}[Condition in terms of a biased distribution on \(U\)]
\label{rem:condition-on-u}
Since \(T\) and \(B\) are functions of \(U\), the condition \(B\perp T\) can be expressed directly in terms of the law of \(U\). Writing \(p_t:=\mathbb{P}(U=t)\) and \(q_t:=\mathbb{P}(U=t+n)\) for \(t\in\{0,\dots,n-1\}\), one has
\[
\mathbb{P}(B=1\mid T=t)=\frac{q_t}{p_t+q_t}.
\]
Thus \(B\perp T\) is equivalent to the ratio \(q_t/(p_t+q_t)\) being constant in \(t\). A common special case is the balanced condition \(p_t=q_t\) for all \(t\), which forces \(B\sim\mathrm{Bernoulli}(1/2)\) and \(B\perp T\).
\end{remark}

\begin{openquestion}
\label{oq:continuous-info}
It remains to be shown how to formulate an analogous information-theoretic statement for a continuous model \((U,V)\in T^2\) with an observation \(K=p(U,V)\in K\), in a way that accounts for measurement precision (e.g.\ via quantization) and avoids pathologies of differential entropy for deterministic mappings.
\end{openquestion}

\section{A topology interlude: bundles over the circle and homotopy lifting}
\label{sec:topology}

This section records standard topology facts that motivate the terms ``fiber bundle'', ``fibration'', and ``homotopy lifting property'' in the examples discussed above. No security interpretation is intended here.

\subsection{The Möbius band as an interval bundle over \(S^1\)}
\label{sec:mobius-to-circle}

Let \(I=[0,1]\). The Möbius band \(M\) is the mapping torus of the reflection \(r:I\to I\), \(r(t)=1-t\):
\[
M := (I\times [0,1])/\!\sim,
\qquad
(t,0)\sim(1-t,1).
\]
Concretely, one takes a square \(I\times[0,1]\) and identifies the two vertical edges with a flip: the point \((t,0)\) on the left edge is glued to \((1-t,1)\) on the right edge.

The projection
\[
\pi_M:M\to S^1,\qquad \pi_M([t,s])=e^{2\pi i s},
\]
makes \(M\) into a (nontrivial) \(I\)-bundle over \(S^1\). The nontriviality is encoded by the monodromy: traversing once around the base circle (as \(s\) goes from \(0\) to \(1\)) identifies the fiber at \(s=0\) with the fiber at \(s=1\) via \(t\mapsto 1-t\).

\subsection{The Klein bottle as a circle bundle over \(S^1\)}
\label{sec:klein-to-circle}

Let \(r:S^1\to S^1\) be a reflection (an orientation-reversing involution). The Klein bottle \(K\) can be presented as the mapping torus of \(r\):
\[
K \cong (S^1\times[0,1])/\!\sim,
\qquad
(x,0)\sim(r(x),1).
\]
The projection
\[
\pi_K:K\to S^1,\qquad \pi_K([x,t]) = e^{2\pi i t},
\]
makes \(K\) into a (nontrivial) \(S^1\)-bundle over \(S^1\). As in the Möbius case, the nontriviality is encoded by monodromy: a loop around the base applies \(r\) to the fiber circle.

\subsection{Homotopy lifting property (HLP) and the lifting diagram}
\label{sec:hlp}

The standard definition of ``fibration'' in homotopy theory is phrased via the homotopy lifting property.

\begin{definition}[Homotopy lifting property]
\label{def:hlp}
A map \(p:E\to B\) is said to have the \emph{homotopy lifting property} (HLP) if for every space \(X\), every map \(\widetilde f:X\to E\), and every homotopy \(H:X\times I\to B\) satisfying \(p\circ \widetilde f = H\circ i_0\) (where \(i_0(x)=(x,0)\)), there exists a map \(\widetilde H:X\times I\to E\) such that
\[
\widetilde H\circ i_0=\widetilde f,
\qquad
p\circ \widetilde H = H.
\]
\end{definition}

\begin{remark}[Commutative diagram form]
\label{rem:hlp-diagram}
\Cref{def:hlp} is equivalently expressed by the existence of a dotted lift making the following diagram commute:
\[
\begin{array}{ccc}
X & \xrightarrow{\ \widetilde f\ } & E \\
\Big\downarrow {i_0} & \nearrow_{\ \exists\,\widetilde H\ } & \Big\downarrow {p} \\
X\times I & \xrightarrow{\ H\ } & B.
\end{array}
\]
\end{remark}

\begin{proposition}[Bundles are fibrations]
\label{prop:bundle-fibration}
If \(p:E\to B\) is a (locally trivial) fiber bundle, then \(p\) satisfies the homotopy lifting property.
\end{proposition}

\begin{proof}[Proof sketch]
By local triviality, each point of \(B\) has a neighborhood \(U\) with \(p^{-1}(U)\cong U\times F\). Over such a neighborhood, homotopies lift by projection to the \(U\)-factor and identity on the \(F\)-factor. A global lift is constructed by covering \(X\times I\) with preimages of trivializing neighborhoods, using compactness (in the \(I\)-direction) to patch local lifts together. Continuity of the resulting lift follows from the gluing lemma.
\end{proof}

\begin{remark}
\label{rem:mobius-klein-hlp}
The projections \(\pi_M:M\to S^1\) and \(\pi_K:K\to S^1\) in \cref{sec:mobius-to-circle,sec:klein-to-circle} are fiber bundles, hence satisfy HLP by \cref{prop:bundle-fibration}. Their nontriviality appears in the monodromy (how a loop in the base acts on the fiber), not as a failure of the lifting property.
\end{remark}

\subsection{Example: homotopy groups of the Klein bottle from the fibration \(S^1\to K\to S^1\)}
\label{sec:klein-les}

This subsection elaborates the relationship between the fibration \(\pi_K:K\to S^1\) and the fundamental group of \(K\). It also explains, in a concrete instance, why HLP (or, more generally, the fibration property) is valuable: it provides a long exact sequence in homotopy.

\begin{lemma}[Higher homotopy groups of a covering space]
\label{lem:covering-higher-pi}
Let \(q:\widetilde X\to X\) be a covering map of path-connected spaces, and fix basepoints \(\tilde x_0\in \widetilde X\), \(x_0=q(\tilde x_0)\). Then for each \(n\ge 2\), the induced map
\[
q_*:\pi_n(\widetilde X,\tilde x_0)\to \pi_n(X,x_0)
\]
is an isomorphism.
\end{lemma}

\begin{proof}
For \(n\ge 2\), the sphere \(S^n\) is simply connected. Any based map \(f:(S^n,*)\to (X,x_0)\) admits a unique based lift \(\tilde f:(S^n,*)\to (\widetilde X,\tilde x_0)\) through \(q\), and based homotopies between such maps lift uniquely as well. These lifting and uniqueness properties imply that \(q_*\) is bijective.
\end{proof}

\begin{corollary}
\label{cor:pi-n-s1}
For all \(n\ge 2\), one has \(\pi_n(S^1)=0\).
\end{corollary}

\begin{proof}
The universal cover \(q:\mathbb{R}\to S^1\) is a covering map. By \cref{lem:covering-higher-pi}, \(q_*:\pi_n(\mathbb{R})\to \pi_n(S^1)\) is an isomorphism for \(n\ge 2\). Since \(\mathbb{R}\) is contractible, \(\pi_n(\mathbb{R})=0\), hence \(\pi_n(S^1)=0\).
\end{proof}

\begin{proposition}[Vanishing of \(\pi_n(K)\) for \(n\ge 2\)]
\label{prop:pi-n-klein}
Let \(\pi_K:K\to S^1\) be the circle bundle in \cref{sec:klein-to-circle}, with fiber \(F\cong S^1\). Then \(\pi_n(K)=0\) for all \(n\ge 2\).
\end{proposition}

\begin{proof}
The long exact sequence in homotopy associated to the fibration \(F\to K\to S^1\) contains, for each \(n\ge 3\), the segment
\[
\pi_n(F)\to \pi_n(K)\to \pi_n(S^1).
\]
By \cref{cor:pi-n-s1}, \(\pi_n(S^1)=0\) for \(n\ge 2\). Also \(F\cong S^1\), hence \(\pi_n(F)=\pi_n(S^1)=0\) for \(n\ge 2\). Therefore \(\pi_n(K)=0\) for \(n\ge 3\).

For \(n=2\), the sequence contains \(\pi_2(F)\to \pi_2(K)\to \pi_2(S^1)\), and both end groups vanish by \cref{cor:pi-n-s1}, so \(\pi_2(K)=0\) as well.
\end{proof}

\begin{proposition}[A short exact sequence for \(\pi_1(K)\)]
\label{prop:pi1-short-exact}
With \(F\cong S^1\) the fiber of \(\pi_K:K\to S^1\), there is a short exact sequence
\[
0 \longrightarrow \pi_1(F) \xrightarrow{i_*} \pi_1(K) \xrightarrow{(\pi_K)_*} \pi_1(S^1) \longrightarrow 0.
\]
In particular, \(\pi_1(K)\) is an extension of \(\mathbb{Z}\) by \(\mathbb{Z}\).
\end{proposition}

\begin{proof}
The long exact sequence contains
\[
\pi_2(S^1)\to \pi_1(F)\xrightarrow{i_*}\pi_1(K)\xrightarrow{(\pi_K)_*}\pi_1(S^1)\to \pi_0(F).
\]
The left term vanishes: \(\pi_2(S^1)=0\) by \cref{cor:pi-n-s1}. The right term \(\pi_0(F)\) is a singleton (equivalently, the trivial pointed set) since \(F\cong S^1\) is path-connected. Exactness at \(\pi_1(F)\) and \(\pi_1(S^1)\) therefore yields the short exact sequence. Finally, \(\pi_1(F)\cong\pi_1(S^1)\cong \mathbb{Z}\).
\end{proof}

\begin{proposition}[A presentation of \(\pi_1(K)\) via monodromy]
\label{prop:pi1-klein-presentation}
Let \(a\in \pi_1(K)\) be represented by a loop that projects under \(\pi_K\) to the generator of \(\pi_1(S^1)\), and let \(b\in \pi_1(K)\) be represented by the generator of the fiber \(\pi_1(F)\). Then
\[
\pi_1(K)\cong \langle a,b \mid a b a^{-1} = b^{-1}\rangle \cong \mathbb{Z}\rtimes_{(-1)}\mathbb{Z},
\]
where the base generator acts on the fiber generator by inversion.
\end{proposition}

\begin{proof}
In the mapping-torus model of \cref{sec:klein-to-circle}, traversing once around the base circle identifies the fiber \(S^1\) with itself via the reflection \(r:S^1\to S^1\). On \(\pi_1(S^1)\cong \mathbb{Z}\), the induced automorphism \(r_*\) is multiplication by \(-1\), since \(r\) reverses orientation on the circle. Equivalently, conjugation by a lift of the base generator sends the fiber generator to its inverse, yielding the relation \(aba^{-1}=b^{-1}\).

The short exact sequence of \cref{prop:pi1-short-exact} together with this action identifies \(\pi_1(K)\) as the semidirect product \(\mathbb{Z}\rtimes_{(-1)}\mathbb{Z}\), which has the stated presentation.
\end{proof}

\begin{remark}[Comparison: \(\pi_1\) of the Möbius band]
\label{rem:pi1-mobius}
For the Möbius projection \(\pi_M:M\to S^1\), the fiber \(I\) is contractible, so \(\pi_1(I)=0\). The long exact sequence therefore yields \(\pi_1(M)\cong \pi_1(S^1)\cong \mathbb{Z}\).
\end{remark}

\section{Conclusion}
\label{sec:conclusion}

The linear example \(f(x,y)=ax+cy\) and the topological example \(p:T^2\to K\) illustrate two distinct mechanisms for non-injectivity:
\begin{itemize}
  \item dimension drop (\(\mathbb{R}^2\to\mathbb{R}\)) yielding positive-dimensional fibers;
  \item quotient/covering behavior (\(T^2\to K\)) yielding finite fibers without changing dimension.
\end{itemize}
For protocol design, fiber geometry is a useful intuition, but distributional analysis (e.g.\ conditional entropy or mutual information) is essential for any substantive hiding claim.

% No external references are required for this self-contained note.

\appendix
\section{Dictionary}
\label{sec:dictionary}

This appendix records a small dictionary of recurring terms. Each entry contains an informal description (what it is, how it is used, why it matters) and a formal definition, followed by an example.

\subsection{Path}
\label{sec:dict-path}

\paragraph{Informal.}
A \emph{path} in a topological space \(B\) is a continuous motion in \(B\) over a unit time interval. Paths are used to compare points (via path-connectedness), to define loop-based invariants such as \(\pi_1(B)\), and to express ``dynamics'' of lifted trajectories in coverings and fibrations. In the present note, a path \(\gamma:I\to K\) in the Klein bottle can be \emph{unrolled} to a lifted path \(\tilde\gamma:I\to T^2\) once a starting sheet is fixed; see \cref{def:uplp,rem:uplp-diagram}.

\begin{definition}[Path and loop]
\label{def:path}
Let \(B\) be a topological space and \(I=[0,1]\). A \emph{path} in \(B\) is a continuous map \(\gamma:I\to B\). If a basepoint \(b_0\in B\) is fixed and \(\gamma(0)=\gamma(1)=b_0\), then \(\gamma\) is called a \emph{(based) loop} at \(b_0\).
\end{definition}

\begin{example}
\label{ex:path-circle}
On \(S^1\subset \mathbb{C}\), the map \(\gamma(t)=e^{2\pi i t}\) is a loop that winds once around the circle. In the bundle \(\pi_K:K\to S^1\) of \cref{sec:klein-to-circle}, this loop in the base is exactly the one whose monodromy is the reflection \(r\) on the fiber; see \cref{prop:pi1-klein-presentation}.
\end{example}

\subsection{Covering map}
\label{sec:dict-covering}

\paragraph{Informal.}
A \emph{covering map} \(p:E\to B\) is a surjective map that is locally a disjoint union of homeomorphisms: every point of \(B\) has a neighborhood \(U\) whose preimage looks like several ``sheets'' each mapping homeomorphically to \(U\). Covering maps are used to study spaces by passing to simpler ones (e.g.\ universal covers) and to formalize ``unrolling'' phenomena: paths and homotopies in the base lift uniquely once a starting sheet is chosen. In the present note, the \(2\)-sheeted picture \(T^2\to K\) is used as a model where non-injectivity comes from a discrete identification rather than dimension drop; see \cref{sec:explicit-map,sec:hiding-heuristic}.

\begin{definition}[Covering map]
\label{def:covering-map}
A map \(p:E\to B\) is a \emph{covering map} if it is surjective and for every \(b\in B\) there exists an open neighborhood \(U\subseteq B\) of \(b\) such that
\[
p^{-1}(U)=\bigsqcup_{\alpha\in A} U_\alpha
\]
as a disjoint union of open sets, and for each \(\alpha\) the restriction \(p|_{U_\alpha}:U_\alpha\to U\) is a homeomorphism.
\end{definition}

\begin{example}
\label{ex:covering-circle}
The exponential map \(q:\mathbb{R}\to S^1\), \(q(t)=e^{2\pi i t}\), is a covering map. This covering is used in \cref{cor:pi-n-s1} to conclude that \(\pi_n(S^1)=0\) for \(n\ge 2\), since higher homotopy groups are invariant under covering maps by \cref{lem:covering-higher-pi}.
\end{example}

\begin{example}
\label{ex:covering-torus-klein}
The standard quotient map \(p:T^2\to K\) is a \(2\)-sheeted covering map (each point of \(K\) has a neighborhood whose preimage is two disjoint copies). In this situation, a base path \(\gamma\) has two possible lifts, corresponding to the two choices of initial point in the fiber; once the initial point is fixed, the lift is unique (\cref{def:uplp,prop:covering-uplp}).
\end{example}

\subsection{Monodromy}
\label{sec:dict-monodromy}

\paragraph{Informal.}
\emph{Monodromy} describes how fibers are ``carried around'' when one traverses a loop in the base of a fibration or bundle. It is the mechanism by which a bundle can be nontrivial even when all fibers look identical. In this note, monodromy is what distinguishes the Klein bottle \(S^1\)-bundle over \(S^1\) from the torus \(S^1\times S^1\): going once around the base acts on the fiber by a reflection, which induces the relation \(aba^{-1}=b^{-1}\) in \(\pi_1(K)\); see \cref{prop:pi1-klein-presentation}.

\begin{definition}[Monodromy (bundle over \(S^1\))]
\label{def:monodromy}
Let \(F\hookrightarrow E\xrightarrow{p} S^1\) be a fiber bundle, and fix a basepoint \(*\in S^1\) with fiber \(F=p^{-1}(*)\). Choosing a loop \(\gamma\) that generates \(\pi_1(S^1,*)\) determines (up to homotopy) a fiber self-map \(m:F\to F\) obtained by lifting \(\gamma\) and comparing the resulting identification of the initial and terminal fibers. The homotopy class of \(m\) is called the \emph{monodromy} of the bundle.
\end{definition}

\begin{example}
\label{ex:monodromy-klein}
For the Klein bundle \(\pi_K:K\to S^1\) presented as the mapping torus of a reflection \(r:S^1\to S^1\), the monodromy is precisely \(r\). On \(\pi_1(S^1)\cong \mathbb{Z}\), this induces multiplication by \(-1\), yielding \(\pi_1(K)\cong \langle a,b\mid aba^{-1}=b^{-1}\rangle\) as in \cref{prop:pi1-klein-presentation}.
\end{example}

\begin{example}
\label{ex:monodromy-mobius}
For the Möbius bundle \(\pi_M:M\to S^1\) with fiber \(I=[0,1]\), the monodromy is the reflection \(t\mapsto 1-t\) of the interval. This is the ``twist'' that makes the bundle nontrivial, even though each fiber is homeomorphic to \(I\).
\end{example}

\section{Further examples of fiber ambiguity}
\label{sec:appendix-examples}

This appendix collects standard examples in which a publicly observed value \(b\in B\) has multiple compatible ``hidden'' preimages in a space \(E\) under a map \(p:E\to B\). The purpose is to emphasize that the phenomenon ``\(\#p^{-1}(b)\ge 2\)'' occurs broadly in topology (not only for \(T^2\to K\)), while reiterating that information-theoretic hiding depends on distributional assumptions (as analyzed in \cref{sec:info-theory,sec:biased-u}).

\subsection{Finite-sheeted coverings as constant-size ambiguity}
\label{sec:appendix-finite-coverings}

Covering maps give the cleanest setting for constant-size ambiguity: when \(p:E\to B\) is a connected \(d\)-sheeted covering, every fiber \(p^{-1}(b)\) has exactly \(d\) points, and paths in \(B\) lift uniquely once an initial lift is chosen (\cref{prop:covering-uplp}).

\begin{example}[\(n\)-fold covering of the circle]
\label{ex:nfold-circle}
Let \(p:S^1\to S^1\) be \(p(z)=z^n\) for an integer \(n\ge 2\). Then \(p\) is an \(n\)-sheeted covering map: each output \(w\in S^1\) has exactly \(n\) preimages, the \(n\)th roots of \(w\).
\end{example}

\begin{example}[Multiplication maps on the torus]
\label{ex:nfold-torus}
Identify \(T^2\cong (\mathbb{R}/\mathbb{Z})^2\). For \(n\ge 2\), the map
\[
p_n:T^2\to T^2,\qquad p_n([u],[v])=([nu],[v])
\]
is an \(n\)-sheeted covering map; similarly \(([u],[v])\mapsto([nu],[nv])\) is an \(n^2\)-sheeted covering.
\end{example}

\subsection{Quotients by finite group actions}
\label{sec:appendix-finite-group-quotients}

The map \(T^2\to K\) can be viewed abstractly as a quotient by a free \(\mathbb{Z}_2\)-action. This viewpoint generalizes uniformly: a free action of a finite group produces a finite-sheeted covering whose fibers are the corresponding orbits.

\begin{definition}[Group action and free action]
\label{def:free-action}
Let \(G\) be a group and \(E\) a topological space. A (left) \emph{action} of \(G\) on \(E\) is a map \(G\times E\to E\), \((g,e)\mapsto g\cdot e\), such that \(1\cdot e=e\) and \((g_1g_2)\cdot e=g_1\cdot(g_2\cdot e)\) for all \(g_1,g_2\in G\) and \(e\in E\). The action is \emph{free} if \(g\cdot e=e\) implies \(g=1\).
\end{definition}

\begin{proposition}[Free finite actions give finite-sheeted coverings]
\label{prop:finite-free-quotient-cover}
Let \(G\) be a finite group acting freely on a Hausdorff topological space \(E\). Let \(B:=E/G\) be the orbit space with the quotient topology, and let \(p:E\to B\) be the quotient map \(e\mapsto [e]\). Then \(p\) is a \(|G|\)-sheeted covering map. In particular, each fiber is an orbit:
\[
p^{-1}([e])=\{g\cdot e:\ g\in G\},
\qquad
\#p^{-1}([e])=|G|.
\]
\end{proposition}

\begin{proof}[Proof sketch]
Fix \(e\in E\). Since the action is free and \(G\) is finite, the points \(\{g\cdot e:\ g\in G\}\) are distinct. Using Hausdorffness, choose pairwise disjoint neighborhoods \(U_g\) of \(g\cdot e\). Let \(U:=p(U_1)\subseteq B\), where \(U_1\) is the neighborhood of \(e\). By shrinking \(U_1\) if necessary, one may ensure that \(g\cdot U_1\subseteq U_g\) for each \(g\in G\). Then
\[
p^{-1}(U)=\bigsqcup_{g\in G} g\cdot U_1,
\]
and each restriction \(p|_{g\cdot U_1}:g\cdot U_1\to U\) is a homeomorphism. Thus \(U\) is evenly covered, proving that \(p\) is a covering map of degree \(|G|\).
\end{proof}

\begin{example}[Antipodal quotient \(S^n\to \mathbb{RP}^n\)]
\label{ex:sphere-to-projective}
Let \(G=\mathbb{Z}_2\) act on \(S^n\) by the antipodal map \(x\mapsto -x\). The action is free, and the quotient \(S^n/G\) is real projective space \(\mathbb{RP}^n\). The quotient map \(S^n\to \mathbb{RP}^n\) is therefore a 2-sheeted covering map.
\end{example}

\begin{example}[Revisiting \(T^2\to K\) as a \(\mathbb{Z}_2\)-quotient]
\label{ex:torus-klein-as-quotient}
In \cref{sec:explicit-map}, the involution \(\tau(u,v)=(u+\tfrac12,1-v)\) defines a free \(\mathbb{Z}_2\)-action on \(T^2\). The quotient \(T^2/\langle\tau\rangle\) is homeomorphic to the Klein bottle \(K\), and the quotient map is the 2-sheeted covering \(p:T^2\to K\) used throughout the note.
\end{example}

\end{document}

